% Chapter 5 -- Batch Files.
% Sources (Law 1): os/milton/batch.c (batch_classify, batch_expand %0..%9,
% %VAR%, %%, batch_eval_if); os/milton/command.c (run_batch: ECHO on at
% entry, @, PAUSE MSG-DOS-0016, CALL; IF hands its command to the command
% dispatcher -- GOTO after IF is not recognized, bead needed; command_repl:
% A:\AUTOEXEC.BAT); pubs/screens/disk/*.BAT (A:\SAMPLES of the practice
% diskette).  COPY CON echoes CTRL-Z as a raw glyph that the text decoder
% cannot read, so no screen of it is shown (bead needed, WL-0119).
% GATES: test-batch test-batch-exec test-autoexec test-j1lu-control-emu test-pubs-screens
\Gates{test-batch, test-batch-exec, test-autoexec, test-j1lu-control-emu, test-pubs-screens}
\chapter{Batch Files}

\section{What a Batch File Does}
\index{batch files}
If you type the same commands often, you can put them in a file and carry
them all out by typing one name.  Such a file is a \term{batch file}.  Its
name has the extension \file{.BAT}, and each line holds one command.  When
you type the name of the batch file, without the extension, InitechDOS
carries out the commands in it, in order, as if you had typed them.

\section{Creating a Batch File}
A batch file is an ordinary text file.  You can make one with any program
that saves plain text, or type it directly with \cmd{COPY CON}, which copies
from the keyboard (\cmd{CON}) to a file:
\begin{Steps}
\item Type \type{COPY CON TODAY.BAT} and press \key{Enter}.
\item Type \type{DATE} and press \key{Enter}.
\item Type \type{DIR} and press \key{Enter}.
\item Press \key{Ctrl-Z}, then \key{Enter}.  InitechDOS displays
  \texttt{1 file(s) copied}.
\end{Steps}
Now type \type{TODAY} and press \key{Enter}: InitechDOS displays the date
and then lists the current directory.  If you make a mistake while typing a
line, correct it with \key{Backspace} before you press \key{Enter}; to
change a line already entered, type the whole file again.

\section{The AUTOEXEC.BAT File}
\index{AUTOEXEC.BAT@\texttt{AUTOEXEC.BAT}}
One batch file is special.  When InitechDOS starts, the command processor
looks for a file named \file{AUTOEXEC.BAT} in the root directory of drive~A
and, if it is there, carries it out before it displays the first prompt.
Put in it the commands you want carried out every time you start the
system.  The practice diskette has one:
\begin{Steps}
\item Type \type{CLS} and press \key{Enter}.
\item Type \type{TYPE AUTOEXEC.BAT} and press \key{Enter}.
\item Type \type{SET} and press \key{Enter}.
\end{Steps}

\Screen[rows=9, label=fig:autoexec]{screens/autoexec.txt}{The practice diskette's AUTOEXEC.BAT}

The first line, \texttt{@ECHO OFF}, stops InitechDOS from displaying the
commands of the batch file as it carries them out.  The second puts the
directory \file{A:\bs APPS} in the \texttt{PATH} environment variable, so
that InitechDOS can find the utility programs from any directory.
\cmd{SET} shows the result (Chapter~6).

\section{Replaceable Parameters}
\index{replaceable parameters}
A batch file can be written to work on whatever you name when you run it.
The words you type after its name are its \term{parameters}.  Inside the
file, \texttt{\%1} stands for the first parameter, \texttt{\%2} for the
second, and so on up to \texttt{\%9}.  The batch files used in the rest of
this chapter are in the directory \file{SAMPLES}:
\begin{Steps}
\item Type \type{CLS} and press \key{Enter}.
\item Type \type{CD SAMPLES} and press \key{Enter}.
\item Type \type{TYPE GREET.BAT} and press \key{Enter} to see the file.
\item Type \type{GREET Peter Friday} and press \key{Enter}.
\end{Steps}

\Screen[rows=11, label=fig:greet]{screens/greet.txt}{A batch file with two parameters}

\texttt{\%1} was replaced by \texttt{Peter} and \texttt{\%2} by
\texttt{Friday}.  \cmd{REM} lines are remarks: InitechDOS does not carry
them out.  \cmd{ECHO} followed by text displays the text.

\section{Making Decisions}
\index{IF@\texttt{IF}}
The \cmd{IF} subcommand carries out a command only when a condition is
true.  \file{CHECK.BAT} uses \cmd{IF EXIST} to say whether the file named
by its parameter is on the disk:
\begin{Steps}
\item Type \type{CLS} and press \key{Enter}, then \type{CD SAMPLES}.
\item Type \type{TYPE CHECK.BAT} and press \key{Enter}.
\item Type \type{CHECK GREET.BAT} and press \key{Enter}.
\item Type \type{CHECK NOFILE.TXT} and press \key{Enter}.
\end{Steps}

\Screen[rows=11, label=fig:check]{screens/check.txt}{Testing whether a file exists}

\cmd{IF NOT} reverses a condition.  \cmd{IF} can also compare two strings,
and test the exit code a program returned (\cmd{IF ERRORLEVEL}); the
\doctitle{InitechDOS Reference} gives the details.
\draft{bead needed: the shipped IF passes its command to the command
processor, so IF ... GOTO is not recognized; DOS 3.3 carries it out}
With \cmd{GOTO} and a label line (a line beginning with a colon), a batch
file can jump to another part of itself when a condition is true:
\texttt{IF NOT EXIST \%1 GOTO NOFILE}.

\section{Pausing and Calling Other Batch Files}
\index{PAUSE@\texttt{PAUSE}}\index{CALL@\texttt{CALL}}
\cmd{PAUSE} stops a batch file until you press a key, so that you can, for
example, put paper in the printer.  \cmd{CALL} runs another batch file and
then carries on with the next line.  \file{COVER.BAT} does both:
\begin{Steps}
\item Type \type{CLS} and press \key{Enter}, then \type{CD SAMPLES}.
\item Type \type{COVER} and press \key{Enter}.
\item When \texttt{Press any key to continue . . .} appears, press the
  space bar.
\end{Steps}

\Screen[rows=8, label=fig:cover]{screens/cover.txt}{A batch file that pauses and calls another}

\section{Stopping a Batch File}
To stop a batch file before it has finished, press \key{Ctrl-C}.  The batch
file ends and the command prompt is displayed.

\section{Batch Subcommands}
These subcommands are used in batch files.  The \doctitle{InitechDOS
Reference}, Chapter~4, describes each one.
\par\medskip
\begin{tabular}{@{}l@{\quad}p{70mm}@{}}
{\imSideHeadFont Subcommand} & {\imSideHeadFont What it does}\\[2pt]
\cmd{CALL}  & Runs another batch file, then continues\\
\cmd{ECHO}  & Displays a message, or turns the display of commands on or off\\
\cmd{FOR}   & Carries out a command once for each item in a list\\
\cmd{GOTO}  & Continues at a label line\\
\cmd{IF}    & Carries out a command when a condition is true\\
\cmd{PAUSE} & Waits for a key\\
\cmd{REM}   & Holds a remark\\
\cmd{SHIFT} & Moves the parameters one place to the left\\
\end{tabular}
